\chapter{Análisis de requisitos}

Este capítulo recoge el análisis de requisitos del framework. Dado que el
producto no es una aplicación final sino una \emph{herramienta de
construcción}, los requisitos se expresan en términos de las capacidades que
el framework debe ofrecer a los desarrolladores que lo usen y, de forma
indirecta, a los usuarios finales de las aplicaciones construidas con él.

\section{Identificación de \emph{stakeholders}}
Se identificaron los siguientes \gls{stakeholders}:

\begin{itemize}
  \item \textbf{El \gls{itc}} como propietario del framework y principal
  consumidor: necesita una base reutilizable que reduzca el coste de arrancar
  y mantener nuevos proyectos.
  \item \textbf{El equipo de desarrollo} que construirá aplicaciones sobre el
  framework: necesita una API predecible, tipada y bien documentada.
  \item \textbf{El \gls{jbvc}} como cliente del primer proyecto real
  construido con el framework: aporta el caso de uso que valida la genericidad.
  \item \textbf{El usuario final} de las aplicaciones: necesita formularios y
  vistas correctas, validadas y coherentes, aunque nunca interactúa con el
  framework directamente.
\end{itemize}

\section{Requisitos funcionales}
La tabla~\ref{tab:rf} resume los requisitos funcionales priorizados según el
esquema MoSCoW (\emph{Must}, \emph{Should}, \emph{Could}).

\begin{table}[!htbp]
\caption{Requisitos funcionales del framework.}
\label{tab:rf}
\centering
\begin{tabular}{|l|p{8.5cm}|c|}
\hline
\textbf{ID} & \textbf{Descripción} & \textbf{Prioridad} \\ \hline
RF-01 & Toda respuesta de la API debe seguir un envoltorio uniforme con
contenido, recuento y lista de incidencias. & Must \\ \hline
RF-02 & El framework debe ofrecer fábricas que generen los endpoints CRUD de
una entidad sin escribir un router a mano. & Must \\ \hline
RF-03 & Cada entidad debe publicar su \gls{jsonschema} en un endpoint
dedicado. & Must \\ \hline
RF-04 & El backend debe exponer un endpoint de descubrimiento que liste las
entidades disponibles y sus capacidades. & Must \\ \hline
RF-05 & El frontend debe disponer de una capa cliente genérica tipada que
consuma cualquier entidad a partir de su \emph{path}. & Must \\ \hline
RF-06 & El frontend debe generar un formulario completo a partir del esquema
de una entidad, sin HTML específico. & Must \\ \hline
RF-07 & El esquema debe poder enriquecerse con pistas de presentación
(\emph{widget}, opciones, visibilidad condicional). & Should \\ \hline
RF-08 & El backend debe ofrecer un \emph{bundle} con los esquemas de todas las
entidades, versionado para \emph{caching}. & Should \\ \hline
RF-09 & El sistema debe soportar autenticación (Firebase y un modo local de
desarrollo) y autorización basada en roles y ACL. & Should \\ \hline
RF-10 & Las entidades núcleo deben mantener histórico de cambios
(versionado). & Could \\ \hline
\end{tabular}
\end{table}

\section{Requisitos no funcionales}
\begin{description}
  \item[RNF-01 (Reutilización).] El framework no debe acoplar ninguna lógica
  a un dominio concreto; el dominio se introduce mediante modelos del
  consumidor.
  \item[RNF-02 (Portabilidad).] El sistema completo debe desplegarse de forma
  reproducible en Windows, macOS y Linux mediante \gls{docker}.
  \item[RNF-03 (Mantenibilidad).] El modelo de datos debe ser la única fuente
  de verdad; el contrato y la interfaz deben derivarse de él.
  \item[RNF-04 (Fiabilidad).] Las piezas críticas deben estar cubiertas por
  pruebas automatizadas en backend y frontend.
  \item[RNF-05 (Eficiencia).] El contrato debe permitir \emph{caching}
  agresivo de los esquemas mediante validadores condicionales (ETag).
  \item[RNF-06 (Compatibilidad).] La nueva capa cliente debe convivir con el
  servicio de backend heredado sin romperlo.
  \item[RNF-07 (Seguridad).] Las credenciales y la verificación de tokens no
  deben quedar expuestas; el modo de desarrollo local no debe estar activo en
  producción.
\end{description}

\section{Casos de uso}
Aunque el framework es una herramienta para desarrolladores, su valor se
entiende mejor a través de los casos de uso que habilita. Se distinguen dos
actores: el \emph{desarrollador consumidor}, que construye una aplicación
usando el framework, y el \emph{usuario final}, que interactúa con la
aplicación resultante. El Algoritmo~\ref{lst:casos} resume las relaciones.

\begin{lstlisting}[style=docker, caption={Diagrama textual de casos de uso}, label={lst:casos}]
 Desarrollador consumidor
   |-- CU-01 Declarar una entidad (modelo ORM)
   |-- CU-02 Exponerla via fabrica CRUD
   |-- CU-03 Obtener su esquema enriquecido
   |-- CU-04 Descubrir entidades disponibles
   |-- CU-05 Generar formulario desde el path
   `-- CU-06 Desplegar el sistema (docker compose)

 Usuario final (de la app construida)
   |-- CU-07 Crear/editar un registro vía formulario
   |-- CU-08 Listar/filtrar registros
   `-- CU-09 Autenticarse (Firebase / local dev)
\end{lstlisting}

A continuación se detallan los casos de uso más representativos.

\begin{table}[!htbp]
\caption{Caso de uso CU-05: generar un formulario desde el path de entidad.}
\label{tab:cu05}
\centering
\begin{tabular}{|l|p{9cm}|}
\hline
\textbf{Identificador} & CU-05 \\ \hline
\textbf{Actor} & Desarrollador consumidor \\ \hline
\textbf{Precondición} & La entidad está expuesta vía fábrica CRUD y publica
su \texttt{schema.json}. \\ \hline
\textbf{Flujo principal} & 1) El desarrollador coloca
\texttt{<ngt-dynamic-form entityPath="/x">}. 2) El componente pide el esquema
a \texttt{SchemaService}. 3) \texttt{schemaToFormly} traduce el esquema. 4) Se
renderiza el formulario con validadores. \\ \hline
\textbf{Postcondición} & El usuario final dispone de un formulario funcional
sin que se haya escrito HTML específico. \\ \hline
\textbf{Flujo alternativo} & Si el esquema trae \texttt{issues} de tipo
ERROR, \texttt{SchemaService} propaga el error y el componente muestra las
incidencias. \\ \hline
\end{tabular}
\end{table}

\begin{table}[!htbp]
\caption{Caso de uso CU-07: crear o editar un registro vía formulario.}
\label{tab:cu07}
\centering
\begin{tabular}{|l|p{9cm}|}
\hline
\textbf{Identificador} & CU-07 \\ \hline
\textbf{Actor} & Usuario final \\ \hline
\textbf{Precondición} & Existe un formulario generado para la entidad. \\ \hline
\textbf{Flujo principal} & 1) El usuario rellena el formulario. 2) Los
validadores derivados del esquema comprueban los datos. 3) Al enviar, se
invoca \texttt{EntityClient.create()/update()}. 4) El backend valida con
Pydantic y responde con el envoltorio. \\ \hline
\textbf{Postcondición} & El registro queda persistido y versionado; las
incidencias se muestran bajo el formulario. \\ \hline
\end{tabular}
\end{table}

\section{Restricciones}
El proyecto está condicionado por el stack tecnológico ya existente
(Python/FastAPI/SQLAlchemy en backend y Angular en frontend), por el
calendario académico del TFG y por la necesidad de no introducir regresiones
en el código heredado que la empresa ya utiliza en producción.
